% chapter04.tex — Graphs (source: ch04-graphs.pdf).
% Book pp. 56–89, including Appendix 1 Vectors, Appendix 2 The Conic
% Sections, Appendix 3 Polar Coordinates. Slice pages 3–4 are book 59
% then 58 (swapped leaves); the body is in book order.

\spivakchapter{4}{Graphs}

Mention the real numbers to a mathematician and the image of a straight line
will probably form in her mind, quite involuntarily. And most likely she will
neither banish nor too eagerly embrace this mental picture of the real numbers.
``Geometric intuition'' will allow her to interpret statements about numbers in
terms of this picture, and may even suggest methods of proving them. Although
the properties of the real numbers which were studied in Part~I are not greatly
illuminated by a geometric picture, such an interpretation will be a great aid
in Part~II.

You are probably already familiar with the conventional method of considering
the straight line as a picture of the real numbers, i.e., of associating to each
real number a point on a line. To do this (Figure~1) we pick, arbitrarily, a
point which we label~$0$, and a point to the right, which we label~$1$. The
point twice as far to the right is labeled~$2$, the point the same distance
from~$0$ to~$1$, but to the left of~$0$, is labeled~$-1$, etc. With this
arrangement, if $a<b$, then the point corresponding to~$a$ lies to the left of
the point corresponding to~$b$. We can also draw rational numbers, such as
$\frac12$, in the obvious way. It is usually taken for granted that the
irrational numbers also somehow fit into this scheme, so that every real number
can be drawn as a point on the line. We will not make too much fuss about
justifying this assumption, since this method of ``drawing'' numbers is
intended solely as a method of picturing certain abstract ideas, and our proofs
will never rely on these pictures (although we will frequently use a picture to
suggest or help explain a proof). Because this geometric picture plays such a
prominent, albeit inessential role, geometric terminology is frequently
employed when speaking of numbers---thus a number is sometimes called a
\emph{point}, and $\mathbf{R}$ is often called the \emph{real line}.

\begin{center}
\figtex{assets/geometric/ch04-fig01}\\[0.6em]
\figlabel{1}
\end{center}

The number $\abs{a-b}$ has a simple interpretation in terms of this geometric
picture: it is the distance between $a$ and $b$, the length of the line segment
which has $a$ as one end point and $b$ as the other. This means, to choose an
example whose frequent occurrence justifies special consideration, that the set
of numbers $x$ which satisfy $\abs{x-a}<\varepsilon$ may be pictured as the
collection of points whose distance from $a$ is less than $\varepsilon$. This
set of points is the ``interval'' from $a-\varepsilon$ to $a+\varepsilon$,
which may also be described as the points corresponding to numbers $x$ with
$a-\varepsilon<x<a+\varepsilon$ (Figure~2).

\begin{center}
\figtex{assets/geometric/ch04-fig02}\\[0.6em]
\figlabel{2}
\end{center}

Sets of numbers which correspond to intervals arise so frequently that it is
desirable to have special names for them. The set $\{x:a<x<b\}$ is denoted by
$(a,b)$ and called the \textbf{open interval} from $a$ to $b$. This notation
naturally creates some ambiguity, since $(a,b)$ is also used to denote a pair
of numbers, but in context it is always clear (or can easily be made clear)
whether one is talking about a pair or an interval. Note that if $a\geq b$,
then $(a,b)=\emptyset$, the set with no elements; in practice, however, it is
almost always assumed (explicitly if one has been careful, and implicitly
otherwise), that whenever an interval $(a,b)$ is mentioned, the number $a$ is
less than $b$.

The set $\{x:a\leq x\leq b\}$ is denoted by $[a,b]$ and is called the
\textbf{closed interval} from $a$ to $b$. This symbol is usually reserved for
the case $a<b$, but it is sometimes used for $a=b$, also. The usual pictures
for the intervals $(a,b)$ and $[a,b]$ are shown in Figure~3; since no
reasonably accurate picture could ever indicate the difference between the two
intervals, various conventions have been adopted. Figure~3 also shows certain
``infinite'' intervals. The set $\{x:x>a\}$ is denoted by $(a,\infty)$, while
the set $\{x:x\geq a\}$ is denoted by $[a,\infty)$; the sets $(-\infty,a)$ and
$(-\infty,a]$ are defined similarly. At this point a standard warning must be
issued: the symbols $\infty$ and $-\infty$, though usually read ``infinity'' and
``minus infinity,'' are \emph{purely} suggestive; there is no number
``$\infty$'' which satisfies $\infty\geq a$ for all numbers $a$. While the
symbols $\infty$ and $-\infty$ will appear in many contexts, it is always
necessary to define these uses in ways that refer only to numbers. The set
$\mathbf{R}$ of all real numbers is also considered to be an ``interval,'' and
is sometimes denoted by $(-\infty,\infty)$.

\begin{center}
\figtex{assets/geometric/ch04-fig03}\\[0.6em]
\figlabel{3}
\end{center}

Of even greater interest to us than the method of drawing numbers is a method
of drawing pairs of numbers. This procedure, probably also familiar to you,
requires a ``coordinate system,'' two straight lines intersecting at right
angles. To distinguish these straight lines, we call one the \emph{horizontal
axis}, and one the \emph{vertical axis}. (More prosaic terminology, such as the
``first'' and ``second'' axes, is probably preferable from a logical point of
view, but most people hold their books, or at least their blackboards, in the
same way, so that ``horizontal'' and ``vertical'' are more descriptive.) Each
of the two axes could be labeled with real numbers, but we can also label
points on the horizontal axis with pairs $(a,0)$ and points on the vertical
axis with pairs $(0,b)$, so that the intersection of the two axes, the
``origin'' of the coordinate system, is labeled $(0,0)$. Any pair $(a,b)$ can
now be drawn as in Figure~4, lying at the vertex of the rectangle whose other
three vertices are labeled $(0,0)$, $(a,0)$, and $(0,b)$. The numbers $a$ and
$b$ are called the \emph{first} and \emph{second coordinates}, respectively, of
the point determined in this way.

\begin{center}
\figtex{assets/geometric/ch04-fig04}\\[0.6em]
\figlabel{4}
\end{center}

Our real concern, let us recall, is a method of drawing functions. Since a
function is just a collection of pairs of numbers, we can draw a function by
drawing each of the pairs in the function. The drawing obtained in this way is
called the \textbf{graph} of the function. In other words, the graph of $f$
contains all the points corresponding to pairs $(x,f(x))$. Since most
functions contain infinitely many pairs, drawing the graph promises to be a
laborious undertaking, but, in fact, many functions have graphs which are quite
easy to draw.

Not surprisingly, the simplest functions of all, the constant functions
$f(x)=c$, have the simplest graphs. It is easy to see that the graph of the
function $f(x)=c$ is a straight line parallel to the horizontal axis, at
distance $c$ from it (Figure~5).

\begin{center}
\figtex{assets/plots/ch04-fig05}\\[0.6em]
\figlabel{5}
\end{center}

The functions $f(x)=cx$ also have particularly simple graphs---straight lines
through $(0,0)$, as in Figure~6. A proof of this fact is indicated in
Figure~7:

\begin{center}
\figtex{assets/plots/ch04-fig06}\\[0.6em]
\figlabel{6}
\end{center}

Let $x$ be some number not equal to $0$, and let $L$ be the straight line which
passes through the origin $O$, corresponding to $(0,0)$, and through the
point $A$, corresponding to $(x,cx)$. A point $A'$, with first coordinate $y$,
will lie on $L$ when the triangle $A'B'O$ is similar to the triangle $ABO$,
thus when
\[
\frac{A'B'}{OB'}=\frac{AB}{OB}=c;
\]
this is precisely the condition that $A'$ corresponds to the pair $(y,cy)$,
i.e., that $A'$ lies on the graph of $f$. The argument has implicitly assumed
that $c>0$, but the other cases are treated easily enough. The number $c$,
which measures the ratio of the sides of the triangles appearing in the proof,
is called the \emph{slope} of the straight line, and a line parallel to this
line is also said to have slope $c$.

\begin{center}
\figtex{assets/geometric/ch04-fig07}\\[0.6em]
\figlabel{7}
\end{center}

This demonstration has neither been labeled nor treated as a formal proof.
Indeed, a rigorous demonstration would necessitate a digression which we are
not at all prepared to follow. The rigorous proof of \emph{any} statement
connecting geometric and algebraic concepts would first require a real proof
(or a precisely stated assumption) that the points on a straight line
correspond in an exact way to the real numbers. Aside from this, it would be
necessary to develop plane geometry as precisely as we intend to develop the
properties of real numbers. Now the detailed development of plane geometry is
a beautiful subject, but it is by no means a prerequisite for the study of
calculus. We shall use geometric pictures only as an aid to intuition; for our
purposes (and for most of mathematics) it is perfectly satisfactory to
\emph{define} the plane to be the set of all pairs of real numbers, and to
\emph{define} straight lines as certain collections of pairs, including, among
others, the collections $\{(x,cx):x\text{ a real number}\}$. To provide this
artificially constructed geometry with all the structure of geometry studied
in high school, one more definition is required. If $(a,b)$ and $(c,d)$ are two
points in the plane, i.e., pairs of real numbers, we \emph{define} the
\textbf{distance} between $(a,b)$ and $(c,d)$ to be
\[
\sqrt{(a-c)^2+(b-d)^2}.
\]

\begin{center}
\figtex{assets/geometric/ch04-fig08}\\[0.6em]
\figlabel{8}
\end{center}

If the motivation for this definition is not clear, Figure~8 should serve as
adequate explanation---with this definition the Pythagorean theorem has been
built into our geometry.\footnote{The fastidious reader might object to this
definition on the grounds that nonnegative numbers are not yet known to have
square roots. This objection is really unanswerable at the moment---the
definition will just have to be accepted with reservations, until this little
point is settled.}

Reverting once more to our informal geometric picture, it is not hard to see
(Figure~9) that the graph of the function $f(x)=cx+d$ is a straight line with
slope $c$, passing through the point $(0,d)$. For this reason, the functions
$f(x)=cx+d$ are called \textbf{linear functions}. Simple as they are, linear
functions occur frequently, and you should feel comfortable working with them.
The following is a typical problem whose solution should not cause any trouble.
Given two distinct points $(a,b)$ and $(c,d)$, find the linear function $f$
whose graph goes through $(a,b)$ and $(c,d)$. This amounts to saying that
$f(a)=b$ and $f(c)=d$. If $f$ is to be of the form $f(x)=\alpha x+\beta$, then
we must have
\begin{align*}
\alpha a+\beta&=b,\\
\alpha c+\beta&=d;
\end{align*}
therefore $\alpha=(d-b)/(c-a)$ and $\beta=b-[(d-b)/(c-a)]a$, so
\[
f(x)=\frac{d-b}{c-a}x+b-\frac{d-b}{c-a}a=\frac{d-b}{c-a}(x-a)+b,
\]
a formula most easily remembered by using the ``point-slope form'' (see
Problem~6).

\begin{center}
\figtex{assets/plots/ch04-fig09}\\[0.6em]
\figlabel{9}
\end{center}

Of course, this solution is possible only if $a\neq c$; the graphs of linear
functions account only for the straight lines which are not parallel to the
vertical axis. The vertical straight lines are not the graph of \emph{any}
function at all; in fact, the graph of a function can never contain even two
distinct points on the same vertical line. This conclusion is immediate from
the definition of a function---two points on the same vertical line correspond
to pairs of the form $(a,b)$ and $(a,c)$ and, by definition, a function cannot
contain $(a,b)$ and $(a,c)$ if $b\neq c$. Conversely, if a set of points in the
plane has the property that no two points lie on the same vertical line, then
it is surely the graph of a function. Thus, the first two sets in Figure~10 are
not graphs of functions and the last two are; notice that the fourth is the
graph of a function whose domain is not all of $\mathbf{R}$, since some
vertical lines have no points on them at all.

\begin{center}
\figtex{assets/geometric/ch04-fig10}\\[0.6em]
\figlabel{10}
\end{center}

After the linear functions the simplest is perhaps the function $f(x)=x^2$. If
we draw some of the pairs in $f$, i.e., some of the pairs of the form
$(x,x^2)$, we obtain a picture like Figure~11.

\begin{center}
\figtex{assets/geometric/ch04-fig11}\\[0.6em]
\figlabel{11}
\end{center}

It is not hard to convince yourself that all the pairs $(x,x^2)$ lie along a
curve like the one shown in Figure~12; this curve is known as a
\textbf{parabola}.

\begin{center}
\figtex{assets/plots/ch04-fig12}\\[0.6em]
\figlabel{12}
\end{center}

Since a graph is just a drawing on paper, made (in this case) with printer's
ink, the question ``Is this what the graph really looks like?'' is hard to
phrase in any sensible manner. No drawing is ever \emph{really} correct since
the line has thickness. Nevertheless, there are some questions which one
\emph{can} ask: for example, how can you be sure that the graph does not look
like one of the drawings in Figure~13? It is easy to see, and even to prove,
that the graph cannot look like~(a); for if $0<x<y$, then $x^2<y^2$, so the
graph should be higher at $y$ than at $x$, which is not the case in~(a). It is
also easy to see, simply by drawing a very accurate graph, first plotting many
pairs $(x,x^2)$, that the graph cannot have a large ``jump'' as in~(b) or a
``corner'' as in~(c). In order to prove these assertions, however, we first
need to say, in a mathematical way, what it means for a function not to have a
``jump'' or ``corner''; these ideas already involve some of the fundamental
concepts of calculus. Eventually we will be able to define them rigorously,
but meanwhile you may amuse yourself by attempting to define these concepts,
and then examining your definitions critically. Later these definitions may be
compared with the ones mathematicians have agreed upon. If they compare
favorably, you are certainly to be congratulated!

\begin{center}
\figtex{assets/geometric/ch04-fig13}\\[0.6em]
\figlabel{13}
\end{center}

The functions $f(x)=x^n$, for various natural numbers $n$, are sometimes called
\textbf{power functions}. Their graphs are most easily compared as in
Figure~14, by drawing several at once.

\begin{center}
\figtex{assets/plots/ch04-fig14}\\[0.6em]
\figlabel{14}
\end{center}

The power functions are only special cases of polynomial functions, introduced
in the previous chapter. Two particular polynomial functions are graphed in
Figure~15, while Figure~16 is meant to give a general idea of the graph of the
polynomial function
\[
f(x)=a_n x^n+a_{n-1}x^{n-1}+\cdots+a_0,
\]
in the case $a_n>0$.

\begin{center}
\figtex{assets/plots/ch04-fig15}\\[0.6em]
\figlabel{15}
\end{center}

In general, the graph of $f$ will have at most $n-1$ ``peaks'' or ``valleys''
(a ``peak'' is a point like $(x,f(x))$ in Figure~16, while a ``valley'' is a
point like $(y,f(y))$). The number of peaks and valleys may actually be much
smaller (the power functions, for example, have at most one valley). Although
these assertions are easy to make, we will not even contemplate giving proofs
until Part~III (once the powerful methods of Part~III are available, the proofs
will be very easy).

\begin{center}
\figtex{assets/geometric/ch04-fig16}\\[0.6em]
\figlabel{16}
\end{center}

Figure~17 illustrates the graphs of several rational functions. The
\textbf{rational functions} exhibit even greater variety than the polynomial
functions, but their behavior will also be easy to analyze once we can use the
derivative, the basic tool of Part~III.

\begin{center}
\figtex{assets/plots/ch04-fig17}\\[0.6em]
\figlabel{17}
\end{center}

Many interesting graphs can be constructed by ``piecing together'' the graphs
of functions already studied. The graph in Figure~18 is made up entirely of
straight lines. The function $f$ with this graph satisfies
\begin{align*}
f\Bigl(\frac1n\Bigr)&=(-1)^{n+1},\\
f\Bigl(-\frac1n\Bigr)&=(-1)^{n+1},\\
f(x)&=1, \qquad \abs{x}\geq 1,
\end{align*}
and is a linear function on each interval $[1/(n+1),1/n]$ and
$[-1/n,-1/(n+1)]$. (The number $0$ is not in the domain of $f$.) Of course, one
can write out an explicit formula for $f(x)$, when $x$ is in
$[1/(n+1),1/n]$; this is a good exercise in the use of linear functions, and
will also convince you that a picture is worth a thousand words.

\begin{center}
\figtex{assets/plots/ch04-fig18}\\[0.6em]
\figlabel{18}
\end{center}

It is actually possible to define, in a much simpler way, a function which
exhibits this same property of oscillating infinitely often near $0$, by using
the sine function, which we will discuss in detail in Chapter~15. As usual, we
are using radian measure, so an angle of $2\pi$ means an angle ``all the way
around'' a circle, an angle of $\pi$ an angle half way around (or $180^\circ$ in
layman's terms), an angle of $\pi/2$ a right angle, etc.

The graph of the sine function is shown in Figure~19.

\begin{center}
\figtex{assets/plots/ch04-fig19}\\[0.6em]
\figlabel{19}
\end{center}

Now consider the function $f(x)=\sin 1/x$. The graph of $f$ is shown in
Figure~20. To draw this graph it helps to first observe that
\begin{align*}
f(x)&=0 && \text{for }x=\frac1\pi,\frac1{2\pi},\frac1{3\pi},\dots,\\
f(x)&=1 && \text{for }x=\frac1{\frac12\pi},\frac1{\frac12\pi+2\pi},\frac1{\frac12\pi+4\pi},\dots,\\
f(x)&=-1 && \text{for }x=\frac1{\frac32\pi},\frac1{\frac32\pi+2\pi},\frac1{\frac32\pi+4\pi},\dots.
\end{align*}
Notice that when $x$ is large, so that $1/x$ is small, $f(x)$ is also small;
when $x$ is ``large negative,'' that is, when $\abs{x}$ is large for
negative $x$, again $f(x)$ is close to $0$, although $f(x)<0$.

\begin{center}
\figtex{assets/plots/ch04-fig20}\\[0.6em]
\figlabel{20}
\end{center}

An interesting modification of this function is $f(x)=x\sin 1/x$. The graph of
this function is sketched in Figure~21. Since $\sin 1/x$ oscillates infinitely
often near $0$ between $1$ and $-1$, the function $f(x)=x\sin 1/x$ oscillates
infinitely often between $x$ and $-x$. The behavior of the graph for $x$ large
or large negative is harder to analyze. Since $\sin 1/x$ is getting close to
$0$, while $x$ is getting larger and larger, there seems to be no telling what
the product will do. It is possible to decide, but this is another question
that is best deferred to Part~III. The graph of $f(x)=x^2\sin 1/x$ has also
been illustrated (Figure~22).

\begin{center}
\figtex{assets/plots/ch04-fig21}\\[0.6em]
\figlabel{21}
\end{center}
\begin{center}
\figtex{assets/plots/ch04-fig22}\\[0.6em]
\figlabel{22}
\end{center}

For these infinitely oscillating functions, it is clear that the graph cannot
hope to be really ``accurate.'' The best we can do is to show part of it, and
leave out the part near $0$ (which is the interesting part). Actually, it is
easy to find much simpler functions whose graphs cannot be ``accurately''
drawn. The graphs of
\[
f(x)=\begin{cases} x^2, & x<1\\ 2, & x\geq 1\end{cases}
\qquad\text{and}\qquad
g(x)=\begin{cases} x^2, & x\leq 1\\ 2, & x>1\end{cases}
\]
can only be distinguished by some convention similar to that used for open and
closed intervals (Figure~23).

\begin{center}
\figtex{assets/plots/ch04-fig23}\\[0.6em]
\figlabel{23}
\end{center}

Out last example is a function whose graph is spectacularly nondrawable:
\[
f(x)=\begin{cases} 0, & x\text{ irrational}\\ 1, & x\text{ rational.}\end{cases}
\]
The graph of $f$ must contain infinitely many points on the horizontal axis and
also infinitely many points on a line parallel to the horizontal axis, but it
must not contain either of these lines entirely. Figure~24 shows the usual
textbook picture of the graph. To distinguish the two parts of the graph, the
dots are placed closer together on the line corresponding to irrational $x$.
(There is actually a mathematical reason behind this convention, but it depends
on some sophisticated ideas, introduced in Problems~21-5 and~21-6.)

\begin{center}
\figtex{assets/geometric/ch04-fig24}\\[0.6em]
\figlabel{24}
\end{center}

The peculiarities exhibited by some functions are so engrossing that it is easy
to forget some of the simplest, and most important, subsets of the plane, which
are not the graphs of functions. The most important example of all is the
\textbf{circle}. A circle with center $(a,b)$ and radius $r>0$ contains, by
definition, all the points $(x,y)$ whose distance from $(a,b)$ is equal to $r$.
The circle thus consists (Figure~25) of all points $(x,y)$ with
\[
\sqrt{(x-a)^2+(y-b)^2}=r
\]
or
\[
(x-a)^2+(y-b)^2=r^2.
\]

\begin{center}
\figtex{assets/geometric/ch04-fig25}\\[0.6em]
\figlabel{25}
\end{center}

The circle with center $(0,0)$ and radius $1$, often regarded as a sort of
standard copy, is called the \emph{unit circle}.

A close relative of the circle is the \textbf{ellipse}. This is defined as the
set of points, the \emph{sum} of whose distances from two ``focus'' points is a
constant. (When the two foci are the same, we obtain a circle.) If, for
convenience, the focus points are taken to be $(-c,0)$ and $(c,0)$, and the sum
of the distances is taken to be $2a$ (the factor $2$ simplifies some algebra),
then $(x,y)$ is on the ellipse if and only if
\[
\sqrt{\bigl(x-(-c)\bigr)^2+y^2}+\sqrt{(x-c)^2+y^2}=2a
\]
or
\[
\sqrt{(x+c)^2+y^2}=2a-\sqrt{(x-c)^2+y^2}
\]
or
\[
x^2+2cx+c^2+y^2=4a^2-4a\sqrt{(x-c)^2+y^2}+x^2-2cx+c^2+y^2
\]
or
\[
4(cx-a^2)=-4a\sqrt{(x-c)^2+y^2}
\]
or
\[
c^2x^2-2cxa^2+a^4=a^2(x^2-2cx+c^2+y^2)
\]
or
\[
(c^2-a^2)x^2-a^2y^2=a^2(c^2-a^2)
\]
or
\[
\frac{x^2}{a^2}+\frac{y^2}{a^2-c^2}=1.
\]
This is usually written simply
\[
\frac{x^2}{a^2}+\frac{y^2}{b^2}=1,
\]
where $b=\sqrt{a^2-c^2}$ (since we must clearly choose $a>c$, it follows that
$a^2-c^2>0$). A picture of an ellipse is shown in Figure~26. The ellipse
intersects the horizontal axis when $y=0$, so that
\[
\frac{x^2}{a^2}=1,\qquad x=\pm a,
\]
and it intersects the vertical axis when $x=0$, so that
\[
\frac{y^2}{b^2}=1,\qquad y=\pm b.
\]

\begin{center}
\figtex{assets/geometric/ch04-fig26}\\[0.6em]
\figlabel{26}
\end{center}

The \textbf{hyperbola} is defined analogously, except that we require the
\emph{difference} of the two distances to be constant. Choosing the points
$(-c,0)$ and $(c,0)$ once again, and the constant difference as $2a$, we
obtain, as the condition that $(x,y)$ be on the hyperbola,
\[
\sqrt{(x+c)^2+y^2}-\sqrt{(x-c)^2+y^2}=\pm 2a,
\]
which may be simplified to
\[
\frac{x^2}{a^2}+\frac{y^2}{a^2-c^2}=1.
\]
In this case, however, we must clearly choose $c>a$, so that $a^2-c^2<0$. If
$b=\sqrt{c^2-a^2}$, then $(x,y)$ is on the hyperbola if and only if
\[
\frac{x^2}{a^2}-\frac{y^2}{b^2}=1.
\]
The picture is shown in Figure~27. It contains two pieces, because the
difference between the distances of $(x,y)$ from $(-c,0)$ and $(c,0)$ may be
taken in two different orders. The hyperbola intersects the horizontal axis
when $y=0$, so that $x=\pm a$, but it never intersects the vertical axis.

\begin{center}
\figtex{assets/plots/ch04-fig27}\\[0.6em]
\figlabel{27}
\end{center}

It is interesting to compare (Figure~28) the hyperbola with $a=b=\sqrt{2}$ and
the graph of the function $f(x)=1/x$. The drawings look quite similar, and the
two sets are actually identical, except for a rotation through an angle of
$\pi/4$ (Problem~23).

\begin{center}
\figtex{assets/plots/ch04-fig28}\\[0.6em]
\figlabel{28}
\end{center}

Clearly no rotation of the plane will change circles or ellipses into the
graphs of functions. Nevertheless, the study of these important geometric
figures can often be reduced to the study of functions. Ellipses, for example,
are made up of the graphs of two functions,
\[
f(x)=b\sqrt{1-(x^2/a^2)},\qquad -a\leq x\leq a
\]
and
\[
g(x)=-b\sqrt{1-(x^2/a^2)},\qquad -a\leq x\leq a.
\]
Of course, there are many other pairs of functions with this same property.
For example, we can take
\[
f(x)=\begin{cases}
b\sqrt{1-(x^2/a^2)}, & 0<x\leq a\\
-b\sqrt{1-(x^2/a^2)}, & -a\leq x\leq 0
\end{cases}
\]
and
\[
g(x)=\begin{cases}
-b\sqrt{1-(x^2/a^2)}, & 0<x\leq a\\
b\sqrt{1-(x^2/a^2)}, & -a\leq x\leq 0.
\end{cases}
\]
We could also choose
\[
f(x)=\begin{cases}
b\sqrt{1-(x^2/a^2)}, & x\text{ rational},\ -a\leq x\leq a\\
-b\sqrt{1-(x^2/a^2)}, & x\text{ irrational},\ -a\leq x\leq a
\end{cases}
\]
and
\[
g(x)=\begin{cases}
-b\sqrt{1-(x^2/a^2)}, & x\text{ rational},\ -a\leq x\leq a\\
b\sqrt{1-(x^2/a^2)}, & x\text{ irrational},\ -a\leq x\leq a.
\end{cases}
\]
But all these other pairs necessarily involve unreasonable functions which jump
around. A proof, or even a precise statement of this fact, is too difficult at
present. Although you have probably already begun to make a distinction between
those functions with reasonable graphs, and those with unreasonable graphs, you
may find it very difficult to state a reasonable definition of reasonable
functions. A mathematical definition of this concept is by no means easy, and a
great deal of this book may be viewed as successive attempts to impose more and
more conditions that a ``reasonable'' function must satisfy. As we define some
of these conditions, we will take time out to ask if we have really succeeded
in isolating the functions which deserve to be called reasonable. The answer,
unfortunately, will always be ``no,'' or at best, a qualified ``yes.''

\subsection*{PROBLEMS}

\pnum{1.} Indicate on a straight line the set of all $x$ satisfying the
following conditions. Also name each set, using the notation for intervals (in
some cases you will also need the $\cup$ sign).
\begin{romans}
\item $\abs{x-3}<1$.
\item $\abs{x-3}\leq 1$.
\item $\abs{x-a}<\varepsilon$.
\item $\abs{x^2-1}<\frac12$.
\item $\dfrac{1}{1+x^2}\geq\frac12$.
\item $\dfrac{1}{1+x^2}\leq a$ (give an answer in terms of $a$, distinguishing
  various cases).
\item $x^2+1\geq 2$.
\item $(x+1)(x-1)(x-2)>0$.
\end{romans}

\pnum{2.} There is a very useful way of describing the points of the closed
interval $[a,b]$ (where we assume, as usual, that $a<b$).
\begin{letters}
\item First consider the interval $[0,b]$, for $b>0$. Prove that if $x$ is in
  $[0,b]$, then $x=tb$ for some $t$ with $0\leq t\leq 1$. What is the
  significance of the number $t$? What is the mid-point of the interval
  $[0,b]$?
\item Now prove that if $x$ is in $[a,b]$, then $x=(1-t)a+tb$ for some $t$ with
  $0\leq t\leq 1$. \textit{Hint:} This expression can also be written as
  $a+t(b-a)$. What is the midpoint of the interval $[a,b]$? What is the point
  $1/3$ of the way from $a$ to $b$?
\item Prove, conversely, that if $0\leq t\leq 1$, then $(1-t)a+tb$ is in
  $[a,b]$.
\item The points of the \emph{open} interval $(a,b)$ are those of the form
  $(1-t)a+tb$ for $0<t<1$.
\end{letters}

\pnum{3.} Draw the set of all points $(x,y)$ satisfying the following
conditions. (In most cases your picture will be a sizable portion of a plane,
not just a line or curve.)
\begin{romans}
\item $x>y$.
\item $x+a>y+b$.
\item $y<x^2$.
\item $y\leq x^2$.
\item $\abs{x-y}<1$.
\item $\abs{x+y}<1$.
\item $x+y$ is an integer.
\item $\dfrac{1}{x+y}$ is an integer.
\item $(x-1)^2+(y-2)^2<1$.
\item $x^2<y<x^4$.
\end{romans}

\pnum{4.} Draw the set of all points $(x,y)$ satisfying the following
conditions:
\begin{romans}
\item $\abs{x}+\abs{y}=1$.
\item $\abs{x}-\abs{y}=1$.
\item $\abs{x-1}=\abs{y-1}$.
\item $\abs{1-x}=\abs{y-1}$.
\item $x^2+y^2=0$.
\item $xy=0$.
\item $x^2-2x+y^2=4$.
\item $x^2=y^2$.
\end{romans}

\pnum{5.} Draw the set of all points $(x,y)$ satisfying the following
conditions:
\begin{romans}
\item $x=y^2$.
\item $\dfrac{y^2}{a^2}-\dfrac{x^2}{b^2}=1$.
\item $x=\abs{y}$.
\item $x=\sin y$.
\end{romans}
\textit{Hint:} You already know the answers when $x$ and $y$ are interchanged.

\pnum{6.}
\begin{letters}
\item Show that the straight line through $(a,b)$ with slope $m$ is the graph
  of the function $f(x)=m(x-a)+b$. This formula, known as the ``point-slope
  form'' is far more convenient than the equivalent expression
  $f(x)=mx+(b-ma)$; it is immediately clear from the point-slope form that the
  slope is $m$, and that the value of $f$ at $a$ is $b$.
\item For $a\neq c$, show that the straight line through $(a,b)$ and $(c,d)$ is
  the graph of the function
  \[
  f(x)=\frac{d-b}{c-a}(x-a)+b.
  \]
\item When are the graphs of $f(x)=mx+b$ and $g(x)=m'x+b'$ parallel straight
  lines?
\end{letters}

\pnum{7.}
\begin{letters}
\item For any numbers $A$, $B$, and $C$, with $A$ and $B$ not both $0$, show
  that the set of all $(x,y)$ satisfying $Ax+By+C=0$ is a straight line
  (possibly a vertical one). \textit{Hint:} First decide when a vertical
  straight line is described.
\item Show conversely that every straight line, including vertical ones, can be
  described as the set of all $(x,y)$ satisfying $Ax+By+C=0$.
\end{letters}

\pnum{8.}
\begin{letters}
\item Prove that the graphs of the functions
  \begin{align*}
  f(x)&=mx+b,\\
  g(x)&=nx+c,
  \end{align*}
  are perpendicular if $mn=-1$, by computing the squares of the lengths of the
  sides of the triangle in Figure~29. (Why is this special case, where the
  lines intersect at the origin, as good as the general case?)
\item Prove that the two straight lines consisting of all $(x,y)$ satisfying
  the conditions
  \begin{align*}
  Ax+By+C&=0,\\
  A'x+B'y+C'&=0,
  \end{align*}
  are perpendicular if and only if $AA'+BB'=0$.
\end{letters}

\begin{center}
\figtex{assets/geometric/ch04-fig29}\\[0.6em]
\figlabel{29}
\end{center}

\pnum{9.}
\begin{letters}
\item Prove, using Problem~1-19, that
  \[
  \sqrt{(x_1+y_1)^2+(x_2+y_2)^2}\leq\sqrt{x_1^2+x_2^2}+\sqrt{y_1^2+y_2^2}.
  \]
\item Prove that
  \begin{multline*}
  \sqrt{(x_3-x_1)^2+(y_3-y_1)^2}
  \leq\sqrt{(x_2-x_1)^2+(y_2-y_1)^2}\\
  +\sqrt{(x_3-x_2)^2+(y_3-y_2)^2}.
  \end{multline*}
  Interpret this inequality geometrically (it is called the ``triangle
  inequality''). When does strict inequality hold?
\end{letters}

\pnum{10.} Sketch the graphs of the following functions, plotting enough points
to get a good idea of the general appearance. (Part of the problem is to make a
reasonable decision how many is ``enough''; the queries posed below are meant
to show that a little thought will often be more valuable than hundreds of
individual points.)
\begin{romans}
\item $f(x)=x+\dfrac1x$. (What happens for $x$ near $0$, and for large $x$?
  Where does the graph lie in relation to the graph of the identity function?
  Why does it suffice to consider only positive $x$ at first?)
\item $f(x)=x-\dfrac1x$.
\item $f(x)=x^2+\dfrac1{x^2}$.
\item $f(x)=x^2-\dfrac1{x^2}$.
\end{romans}

\pnum{11.} Describe the general features of the graph of $f$ if
\begin{romans}
\item $f$ is even.
\item $f$ is odd.
\item $f$ is nonnegative.
\item $f(x)=f(x+a)$ for all $x$ (a function with this property is called
  \textbf{periodic}, with \textbf{period} $a$.
\end{romans}

\pnum{12.} Graph the functions $f(x)=\sqrt[m]{x}$ for $m=1,2,3,4$. (There is an
easy way to do this, using Figure~14. Be sure to remember, however, that
$\sqrt[m]{x}$ means the \emph{positive} $m$th root of $x$ when $m$ is even; you
should also note that there will be an important difference between the graphs
when $m$ is even and when $m$ is odd.)

\pnum{13.}
\begin{letters}
\item Graph $f(x)=\abs{x}$ and $f(x)=x^2$.
\item Graph $f(x)=\abs{\sin x}$ and $f(x)=\sin^2 x$. (There is an important
  difference between the graphs, which we cannot yet even describe rigorously.
  See if you can discover what it is; part~(a) is meant to be a clue.)
\end{letters}

\pnum{14.} Describe the graph of $g$ in terms of the graph of $f$ if
\begin{romans}
\item $g(x)=f(x)+c$.
\item $g(x)=f(x+c)$. (It is easy to make a mistake here.)
\item $g(x)=cf(x)$. (Distinguish the cases $c=0$, $c>0$, $c<0$.)
\item $g(x)=f(cx)$.
\item $g(x)=f(1/x)$.
\item $g(x)=f(\abs{x})$.
\item $g(x)=\abs{f(x)}$.
\item $g(x)=\max(f,0)$.
\item $g(x)=\min(f,0)$.
\item $g(x)=\max(f,1)$.
\end{romans}

\pnum{15.} Draw the graph of $f(x)=ax^2+bx+c$. \textit{Hint:} Use the methods
of Problem~1-18.

\pnum{16.} Suppose that $A$ and $C$ are not both $0$. Show that the set of all
$(x,y)$ satisfying
\[
Ax^2+Bx+Cy^2+Dy+E=0
\]
is either a parabola, an ellipse, or an hyperbola (or a ``degenerate case'':
two lines [either intersecting or parallel], one line, a point, or
$\emptyset$). \textit{Hint:} The case $C=0$ is essentially Problem~15, and the
case $A=0$ is just a minor variant. Now consider separately the cases where
$A$ and $B$ are both positive or negative, and where one is positive while the
other is negative. When do we have a circle?

\pnum{17.} The symbol $[x]$ denotes the largest integer which is $\leq x$.
Thus, $[2.1]=[2]=2$ and $[-0.9]=[-1]=-1$. Draw the graph of the following
functions (they are all quite interesting, and several will reappear frequently
in other problems).
\begin{romans}
\item $f(x)=[x]$.
\item $f(x)=x-[x]$.
\item $f(x)=\sqrt{x-[x]}$.
\item $f(x)=[x]+\sqrt{x-[x]}$.
\item $f(x)=\bigl[\frac1x\bigr]$.
\item $f(x)=\dfrac{1}{\bigl[\frac1x\bigr]}$.
\end{romans}

\pnum{18.} Graph the following functions.
\begin{romans}
\item $f(x)=\{x\}$, where $\{x\}$ is defined to be the distance from $x$ to the
  nearest integer.
\item $f(x)=\{2x\}$.
\item $f(x)=\{x\}+\frac12\{2x\}$.
\item $f(x)=\{4x\}$.
\item $f(x)=\{x\}+\frac12\{2x\}+\frac14\{4x\}$.
\end{romans}

Many functions may be described in terms of the decimal expansion of a number.
Although we will not be in a position to describe infinite decimals rigorously
until Chapter~23, your intuitive notion of infinite decimals should suffice to
carry you through the following problem, and others which occur before
Chapter~23. There is one ambiguity about infinite decimals which must be
eliminated: Every decimal ending in a string of 9's is equal to another ending
in a string of 0's (e.g., $1.23999\ldots=1.24000\ldots$). We will always use
the one ending in 9's.

\pnum{*19.} Describe as best you can the graphs of the following functions (a
complete picture is usually out of the question).
\begin{romans}
\item $f(x)={}$the 1st number in the decimal expansion of $x$.
\item $f(x)={}$the 2nd number in the decimal expansion of $x$.
\item $f(x)={}$the number of 7's in the decimal expansion of $x$ if this number
  is finite, and $0$ otherwise.
\item $f(x)=0$ if the number of 7's in the decimal expansion of $x$ is finite,
  and $1$ otherwise.
\item $f(x)={}$the number obtained by replacing all digits in the decimal
  expansion of $x$ which come after the first 7 (if any) by $0$.
\item $f(x)=0$ if 1 never appears in the decimal expansion of $x$, and $n$ if 1
  first appears in the $n$th place.
\end{romans}

\pnum{*20.} Let
\[
f(x)=\begin{cases}
0, & x\text{ irrational}\\
\dfrac1q, & x=\dfrac pq\text{ rational in lowest terms.}
\end{cases}
\]
(A number $p/q$ is in \textbf{lowest terms} if $p$ and $q$ are integers with no
common factor, and $q>0$). Draw the graph of $f$ as well as you can (don't
sprinkle points randomly on the paper; consider first the rational numbers with
$q=2$, then those with $q=3$, etc.).

\pnum{21.}
\begin{letters}
\item The points on the graph of $f(x)=x^2$ are the ones of the form
  $(x,x^2)$. Prove that each such point is equidistant from the point
  $\bigl(0,\frac14\bigr)$ and the graph of $g(x)=-\frac14$. (See Figure~30.)
\item Given a horizontal line $L$, the graph of $g(x)=\gamma$, and a point
  $P=(\alpha,\beta)$ not on $L$, so that $\gamma\neq\beta$, show that the set of
  all points $(x,y)$ equidistant from $P$ and $L$ is the graph of a function of
  the form $f(x)=ax^2+bx+c$. What is this set if $\gamma=\beta$?
\end{letters}

\begin{center}
\figtex{assets/plots/ch04-fig30}\\[0.6em]
\figlabel{30}
\end{center}

\pnum{*22.}
\begin{letters}
\item Show that the square of the distance from $(c,d)$ to $(x,mx)$ is
  \[
  x^2(m^2+1)+x(-2md-2c)+d^2+c^2.
  \]
  Using Problem~1-18 to find the minimum of these numbers, show that the
  distance from $(c,d)$ to the graph of $f(x)=mx$ is
  \[
  \frac{\abs{cm-d}}{\sqrt{m^2+1}}.
  \]
\item Find the distance from $(c,d)$ to the graph of $f(x)=mx+b$. (Reduce this
  case to part~(a).)
\end{letters}

\pnum{*23.}
\begin{letters}
\item Using Problem~22, show that the numbers $x'$ and $y'$ indicated in
  Figure~31 are given by
  \begin{align*}
  x'&=\frac1{\sqrt2}\,x+\frac1{\sqrt2}\,y,\\
  y'&=-\frac1{\sqrt2}\,x+\frac1{\sqrt2}\,y.
  \end{align*}
\item Show that the set of all $(x,y)$ with $(x'/\sqrt2)^2-(y'/\sqrt2)^2=1$ is
  the same as the set of all $(x,y)$ with $xy=1$.
\end{letters}

\begin{center}
\figtex{assets/geometric/ch04-fig31}\\[0.6em]
\figlabel{31}
\end{center}

\section*{APPENDIX 1.\quad VECTORS}

Suppose that $v$ is a point in the plane; in other words, $v$ is a pair of
numbers
\[
v=(v_1,v_2).
\]
For convenience, we will use this convention that subscripts indicate the first
and second pairs of a point that has been described by a single letter. Thus,
if we mention the points $w$ and $z$, it will be understood that $w$ is the pair
$(w_1,w_2)$, while $z$ is the pair $(z_1,z_2)$.

Instead of the actual pair of numbers $(v_1,v_2)$, we often picture $v$ as an
arrow from the origin $O$ to this point (Figure~1), and we refer to these
arrows as \emph{vectors} in the plane. Of course, we've haven't really said
anything new yet, we've simply introduced an alternate term for a point of the
plane, and another mental picture. The real point of the new terminology is to
emphasize that we are going to do some new things with points in the plane.

\begin{center}
\figtex{assets/geometric/ch04-fig32}\\[0.6em]
\figlabel{1}
\end{center}

For example, suppose that we have two vectors (i.e., points) in the plane,
\[
v=(v_1,v_2),\qquad w=(w_1,w_2).
\]
Then we can define a new vector (a new point of the plane) $v+w$ by the
equation
\begin{equation}
v+w=(v_1+w_1,\,v_2+w_2).\tag{1}
\end{equation}
Notice that all the letters on the right side of this equation are numbers, and
the $+$ sign is just our usual addition of numbers. On the other hand, the $+$
sign on the left side is new: previously, the sum of two points in the plane
wasn't defined, and we've simply used equation~(1) as a \emph{definition}.

A very fussy mathematician might want to use some new symbol for this newly
defined operation, like
\[
v\mathbin{\boldsymbol{+}}w, \qquad\text{or perhaps}\qquad v\oplus w,
\]
but there's really no need to insist on this; since $v+w$ hasn't been defined
before, there's no possibility of confusion, so we might as well keep the
notation simple.

Of course, any one can make new notation; for example, since it's our
definition, we could just as well have defined $v+w$ as
$(v_1+w_1\cdot w_2,\,v_2+w_1^2)$, or by some other equally weird formula. The
real question is, does our new construction have any particular significance?

Figure~2 shows two vectors $v$ and $w$, as well as the point
\[
(v_1+w_1,\,v_2+w_2),
\]
which, for the moment, we have simply indicated in the usual way, without
drawing an arrow. Note that it is easy to compute the slope of the line $L$
between $v$ and our new point: as indicated in Figure~2, this slope is just
\[
\frac{(v_2+w_2)-v_2}{(v_1+w_1)-v_1}=\frac{w_2}{w_1},
\]
and this, of course, is the slope of our vector $w$, from the origin $O$ to
$(w_1,w_2)$. In other words, the line $L$ is parallel to $w$.

\begin{center}
\figtex{assets/geometric/ch04-fig33}\\[0.6em]
\figlabel{2}
\end{center}

Similarly, the slope of the line $M$ between $(w_1,w_2)$ and our new point is
\[
\frac{(v_2+w_2)-w_2}{(v_1+w_1)-v_2}=\frac{v_2}{v_1},
\]
which is the slope of the vector $v$; so $M$ is parallel to $v$. In short, the
new point $v+w$ lies on the parallelogram having $v$ and $w$ as sides. When we
draw $v+w$ as an arrow (Figure~3), it points along the diagonal of this
parallelogram. In physics, vectors are used to symbolize forces, and the sum of
two vectors represents the resultant force when two different forces are
applied simultaneously to the same object.

\begin{center}
\figtex{assets/geometric/ch04-fig34}\\[0.6em]
\figlabel{3}
\end{center}

Figure~4 shows another way of visualizing the sum $v+w$. If we use ``$w$'' to
denote an arrow parallel to $w$, and having the same length, but starting at
$v$ instead of at the origin, then $v+w$ is the vector from $O$ to the final
endpoint; thus we get to $v+w$ by first following $v$, and then following $w$.

\begin{center}
\figtex{assets/geometric/ch04-fig35}\\[0.6em]
\figlabel{4}
\end{center}

Many of the properties of $+$ for ordinary numbers also hold for this new $+$
for vectors. For example, the ``commutative law''
\[
v+w=w+v,
\]
is obvious from the geometric picture, since the parallelogram spanned by $v$
and $w$ is the same as the parallelogram spanned by $w$ and $v$. It is also
easily checked analytically, since it states that
\[
(v_1+w_1,\,v_2+w_2)=(w_1+v_1,\,w_2+v_2),
\]
and thus simply depends on the commutative law for numbers:
\begin{align*}
v_1+w_1&=w_1+v_1,\\
v_2+w_2&=w_2+v_2.
\end{align*}
Similarly, unraveling definitions, we find the ``associative law''
\[
[v+w]+z=v+[w+z].
\]
Figure~5 indicates a method of finding $v+w+z$.

\begin{center}
\figtex{assets/geometric/ch04-fig36}\\[0.6em]
\figlabel{5}
\end{center}

The origin $O=(0,0)$ is an ``additive identity,''
\[
O+v=v+O=v,
\]
and if we define
\[
-v=(-v_1,-v_2),
\]
then we also have
\[
v+(-v)=-v+v=O.
\]
Naturally we can also define
\[
w-v=w+(-v),
\]
exactly as with numbers; equivalently,
\[
w-v=(w_1-v_1,\,w_2-v_2).
\]

Just as with numbers, our definition of $w-v$ simply means that it satisfies
\[
v+(w-v)=w.
\]
Figure~6(a) shows $v$ and an arrow ``$w-v$'' that is parallel to $w-v$ but that
starts at the endpoint of $v$. As we established with Figure~4, the vector from
the origin to the endpoint of this arrow is just $v+(w-v)=w$ (Figure~6(b)). In
other words, we can picture $w-v$ geometrically as the arrow that goes from $v$
to $w$ (except that it must then be moved back to the origin).

\begin{center}
\figtex{assets/geometric/ch04-fig37}\\[0.6em]
\figlabel{6}
\end{center}

There is also a way of multiplying a number by a vector: For a number $a$ and a
vector $v=(v_1,v_2)$, we define
\[
a\cdot v=(av_1,av_2)
\]
(We sometimes simply write $av$ instead of $a\cdot v$; of course, it is then
especially important to remember that $v$ denotes a vector, rather than a
number.) The vector $a\cdot v$ points in the same direction as $v$ when $a>0$
and in the opposite direction when $a<0$ (Figure~7).

\begin{center}
\figtex{assets/geometric/ch04-fig38}\\[0.6em]
\figlabel{7}
\end{center}

You can easily check the following formulas:
\begin{align*}
a\cdot(b\cdot v)&=(ab)\cdot v,\\
1\cdot v&=v,\\
0\cdot v&=O,\\
-1\cdot v&=-v.
\end{align*}
Notice that we have only defined a product of a number and a vector, we have
not defined a way of `multiplying' two vectors, to get another
vector.\footnote{If you jump to Chapter~25, you'll find that there is an
important way of defining a product, but this is something very special for the
plane---it doesn't work for vectors in 3-space, for example, even though the
other constructions do.} However, there are various ways of `multiplying'
vectors to get numbers, which are explored in the following problems.

\subsection*{PROBLEMS}

\pnum{1.} Given a point $v$ of the plane, let $R_{\theta}(v)$ be the result of
rotating $v$ around the origin through an angle of $\theta$ (Figure~8). The aim
of this problem is to obtain a formula for $R_{\theta}$, with minimal
calculation.
\begin{letters}
\item Show that
  \begin{align*}
  R_{\theta}(1,0)&=(\cos\theta,\sin\theta),
  &&\text{[we should really write $R_{\theta}((1,0))$, etc.]}\\
  R_{\theta}(0,1)&=(-\sin\theta,\cos\theta).
  \end{align*}
\item Explain why we have
  \begin{align*}
  R_{\theta}(v+w)&=R_{\theta}(v)+R_{\theta}(w),\\
  R_{\theta}(a\cdot w)&=a\cdot R_{\theta}(w).
  \end{align*}
\item Now show that for any point $(x,y)$ we have
  \[
  R_{\theta}(x,y)=(x\cos\theta-y\sin\theta,\,x\sin\theta+y\cos\theta).
  \]
\item Use this result to give another solution to Problem~4-23.
\end{letters}

\begin{center}
\figtex{assets/geometric/ch04-fig39}\\[0.6em]
\figlabel{8}
\end{center}

\pnum{2.} Given $v$ and $w$, we define the \emph{number}
\[
v\cdot w=v_1w_1+v_2w_2;
\]
this is often called the `dot product' or `scalar product' of $v$ and $w$
(`scalar' being a rather old-fashioned word for a number, as opposed to a
vector).
\begin{letters}
\item Given $v$, find a vector $w$ such that $v\cdot w=0$. Now describe the set
  of all such vectors $w$.
\item Show that
  \begin{align*}
  v\cdot w&=w\cdot v\\
  v\cdot(w+z)&=v\cdot w+v\cdot z
  \end{align*}
  and that
  \[
  a\cdot(v\cdot w)=(a\cdot v)\cdot w=v\cdot(a\cdot w).
  \]
  Notice that the last of these equations involves \emph{three} products: the
  dot product $\cdot$ of two vectors; the product $\cdot$ of a number and a
  vector; and the ordinary product $\cdot$ of two numbers.
\item Show that $v\cdot v\geq 0$, and that $v\cdot v=0$ only when $v=O$. Hence
  we can define the \emph{norm} $\lVert v\rVert$ as
  \[
  \lVert v\rVert=\sqrt{v\cdot v},
  \]
  which will be $0$ only for $v=O$. What is the geometric interpretation of the
  norm?
\item Prove that
  \[
  \lVert v+w\rVert\leq\lVert v\rVert+\lVert w\rVert,
  \]
  and that equality holds if and only if $v=0$ or $w=0$ or $w=a\cdot v$ for
  some number $a>0$.
\item Show that
  \[
  v\cdot w=\frac{\lVert v+w\rVert^2-\lVert v-w\rVert^2}{4}.
  \]
\end{letters}

\pnum{3.}
\begin{letters}
\item Let $R_{\theta}$ be rotation by an angle of $\theta$ (Problem~1). Show
  that
  \[
  R_{\theta}(v)\cdot R_{\theta}(w)=v\cdot w.
  \]
\item Let $e=(1,0)$ be the vector of length $1$ pointing along the first axis,
  and let $w=(\cos\theta,\sin\theta)$; this is a vector of length $1$ that makes
  an angle of $\theta$ with the first axis (compare Problem~1). Calculate that
  \[
  e\cdot w=\cos\theta.
  \]
  Conclude that in general
  \[
  v\cdot w=\lVert v\rVert\cdot\lVert w\rVert\cdot\cos\theta,
  \]
  where $\theta$ is the angle between $v$ and $w$.
\end{letters}

\pnum{4.} Given two vectors $v$ and $w$, we'd expect to have a simple formula,
involving the coordinates $v_1$, $v_2$, $w_1$, $w_2$, for the area of the
parallelogram they span. Figure~9 indicates a strategy for finding such a
formula: since the triangle with vertices $w$, $A$, $v+w$ is congruent to the
triangle $OBv$, we can reduce the problem to an easier one where one side of
the parallelogram lies along the horizontal axis:
\begin{letters}
\item The line $L$ passes through $v$ and is parallel to $w$, so has slope
  $w_2/w_1$. Conclude that the point $B$ has coordinate
  \[
  \frac{v_1w_2-w_1v_2}{w_2},
  \]
  and that the parallelogram therefore has area
  \[
  \det(v,w)=v_1w_2-w_1v_2.
  \]
\end{letters}
This formula, which defines the \emph{determinant} $\det$, certainly seems to
be simple enough, but it can't really be true that $\det(v,w)$ always gives the
area. After all, we clearly have
\[
\det(w,v)=-\det(v,w),
\]
so sometimes $\det$ will be negative! Indeed, it is easy to see that our
``derivation'' made all sorts of assumptions (that $w_2$ was positive, that $B$
had a positive coordinate, etc.) Nevertheless, it seems likely that
$\det(v,w)$ is $\pm$ the area; the next problem gives an independent proof.

\begin{center}
\figtex{assets/geometric/ch04-fig40}\\[0.6em]
\figlabel{9}
\end{center}

\pnum{5.}
\begin{letters}
\item If $v$ points along the positive horizontal axis, show that $\det(v,w)$
  is the area of the parallelogram spanned by $v$ and $w$ for $w$ above the
  horizontal axis ($w_2>0$), and the negative of the area for $w$ below this
  axis.
\item If $R_{\theta}$ is rotation by an angle of $\theta$ (Problem~1), show that
  \[
  \det(R_{\theta}v,\,R_{\theta}w)=\det(v,w).
  \]
  Conclude that $\det(v,w)$ is the area of the parallelogram spanned by $v$ and
  $w$ when the rotation from $v$ to $w$ is counterclockwise, and the negative
  of the area when it is clockwise.
\end{letters}

\pnum{6.} Show that
\begin{align*}
\det(v,w+z)&=\det(v,w)+\det(v,z)\\
\det(v+w,z)&=\det(v,z)+\det(w,z)
\end{align*}
and that
\[
a\det(v,w)=\det(a\cdot v,w)=\det(v,a\cdot w).
\]

\pnum{7.} Using the method of Problem~3, show that
\[
\det(v,w)=\lVert v\rVert\cdot\lVert w\rVert\cdot\sin\theta,
\]
which is also obvious from the geometric interpretation (Figure~10).

\begin{center}
\figtex{assets/geometric/ch04-fig41}\\[0.6em]
\figlabel{10}
\end{center}

\section*{APPENDIX 2.\quad THE CONIC SECTIONS}

Although we will be concerned almost exclusively with figures in the plane,
defined formally as the set of all pairs of real numbers, in this Appendix we
want to consider three-dimensional space, which we can describe in terms of
triples of real numbers, using a ``three-dimensional coordinate system,''
consisting of three straight lines intersecting at right angles (Figure~1).
Our \emph{horizontal} and \emph{vertical} axes now mutate to two axes in a
horizontal plane, with the third axis perpendicular to both.

\begin{center}
\figtex{assets/geometric/ch04-fig42}\\[0.6em]
\figlabel{1}
\end{center}

One of the simplest subsets of this three-dimensional space is the (infinite)
\emph{cone} illustrated in Figure~2; this cone may be produced by rotating a
``generating line,'' of slope $C$ say, around the third axis.

\begin{center}
\figtex{assets/geometric/ch04-fig43}\\[0.6em]
\figlabel{2}
\end{center}

For any given first two coordinates $x$ and $y$, the point $(x,y,0)$ in the
horizontal plane has distance $\sqrt{x^2+y^2}$ from the origin, and thus
\begin{equation}
(x,y,z)\text{ is on the cone if and only if }z=\pm C\sqrt{x^2+y^2}.\tag{1}
\end{equation}

We can descend from these three-dimensional vistas to the more familiar
two-dimensional one by asking what happens when we intersect this cone with
some plane $P$ (Figure~3).

\begin{center}
\figtex{assets/geometric/ch04-fig44}\\[0.6em]
\figlabel{3}
\end{center}

If the plane is parallel to the horizontal plane, there's certainly no
mystery---the intersection is just a circle. Otherwise, the plane $P$
intersects the horizontal plane in a straight line. We can make things a lot
simpler for ourselves if we rotate everything around the vertical axis so that
this intersection line points straight out from the plane of the paper, while
the first axis is in the usual position that we are familiar with. The plane
$P$ is thus viewed ``straight on,'' so that all we see (Figure~4) is its
intersection $L$ with the plane of the first and third axes; from this
view-point the cone itself simply appears as two straight lines.

\begin{center}
\figtex{assets/geometric/ch04-fig45}\\[0.6em]
\figlabel{4}
\end{center}

If this line $L$ happens to be vertical, consisting of all points $(a,z)$ for
some $a$, then equation~(1) says that the intersection of the cone and the
plane consists of all points $(a,y,z)$ with
\[
z^2-C^2y^2=C^2a^2,
\]
which is an hyperbola.

Otherwise, in the plane of the first and third axes, the line $L$ can be
described as the collection of all points of the form
\[
(x,Mx+B),
\]
where $M$ is the slope of $L$. For an arbitrary point $(x,y,z)$ it follows that
\begin{equation}
(x,y,z)\text{ is in the plane }P\text{ if and only if }z=Mx+B.\tag{2}
\end{equation}

Combining~(1) and~(2), we see that $(x,y,z)$ is in the intersection of the cone
and the plane if and only if
\begin{equation}
Mx+B=\pm C\sqrt{x^2+y^2}.\tag{$*$}
\end{equation}

Now we have to choose coordinate axes in the plane $P$. We can choose $L$ as
the first axis, measuring distances from the intersection $Q$ with the
horizontal plane (Figure~5); for the second axis we just choose the line
through $Q$ parallel to our original second axis. If the first coordinate of a
point in $P$ with respect to these axes is $x$, then the first coordinate of
this point with respect to the original axes can be written in the form
\[
\alpha x+\beta
\]
for some $\alpha$ and $\beta$. On the other hand, if the second coordinate of
the point with respect to these axes is $y$, then $y$ is also the second
coordinate with respect to the original axes.

\begin{center}
\figtex{assets/geometric/ch04-fig46}\\[0.6em]
\figlabel{5}
\end{center}

Consequently, ($*$) says that the point lies on the intersection of the plane
and the cone if and only if
\[
M(\alpha x+\beta)+B=\pm C\sqrt{(\alpha x+\beta)^2+y^2}.
\]
Although this looks fairly complicated, after squaring we can write this as
\[
C^2y^2-\alpha^2(M^2-C^2)x^2+Ex+F=0
\]
for some $E$ and $F$ that we won't bother writing out.

Now Problem~4-16 indicates that this is either a parabola, an ellipse, or an
hyperbola. In fact, looking a little more closely at the solution, we see that
the values of $E$ and $F$ are irrelevant:
\begin{steps}
\item If $M=\pm C$ we obtain a parabola;
\item If $C^2>M^2$ we obtain an ellipse;
\item If $C^2<M^2$ we obtain an hyperbola.
\end{steps}
These analytic conditions are easy to interpret geometrically (Figure~6):
\begin{steps}
\item If our plane is parallel to one of the generating lines of the cone we
  obtain a parabola;
\item If our plane slopes less than the generating line of the cone (so that
  our intersection omits one half of the cone) we obtain an ellipse;
\item If our plane slopes more than the generating line of the cone we obtain
  an hyperbola.
\end{steps}

\begin{center}
\includesvg[width=0.95\textwidth,inkscapelatex=false]{assets/rendered/ch04-fig47}\\[0.6em]
\figlabel{6}
\end{center}

In fact, the very names of these ``conic sections'' are related to this
description. The word \emph{parabola} comes from a Greek root meaning
`alongside,' the same root that appears in parable, not to mention paradigm,
paradox, paragon, paragraph, paralegal, parallax, parallel, and even parachute.
\emph{Ellipse} comes from a Greek root meaning `defect,' or omission, as in
ellipsis (an omission, \ldots\ or the dots that indicate it). And
\emph{hyperbola} comes from a Greek root meaning `throwing beyond,' or excess.
With the currency of words like hyperactive, hypersensitive, and
hyperventilate, not to mention hype, one can probably say, without risk of
hyperbole, that this root is familiar to almost everyone.\footnote{Although the
correspondence between these roots and the geometric picture correspond so
beautifully, for the sake of dull accuracy it has to be reported that the
Greeks originally applied the words to describe features of certain equations
involving the conic sections.}

\subsection*{PROBLEMS}

\pnum{1.} Consider a cylinder with a generator perpendicular to the horizontal
plane (Figure~7); the only requirement for a point $(x,y,z)$ to lie on this
cylinder is that $(x,y)$ lies on a circle:
\[
x^2+y^2=C^2.
\]
Show that the intersection of a plane with this cylinder can be described by
an equation of the form
\[
(\alpha x+\beta)^2+y^2=C^2.
\]
What possibilities are there?

\begin{center}
\includesvg[width=0.42\textwidth,inkscapelatex=false]{assets/rendered/ch04-fig48}\\[0.6em]
\figlabel{7}
\end{center}

\pnum{2.} In Figure~8, the sphere $S_1$ has the same diameter as the cylinder,
so that its equator $C_1$ lies along the cylinder; it is also tangent to the
plane $P$ at $F_1$. Similarly, the equator $C_2$ of $S_2$ lies along the
cylinder, and $S_2$ is tangent to $P$ at $F_2$.
\begin{letters}
\item Let $z$ be any point on the intersection of $P$ and the cylinder. Explain
  why the length of the line from $z$ to $F_1$ is equal to the length of the
  vertical line $L$ from $z$ to $C_1$.
\item By proving a similar fact for the length of the line from $z$ to $F_2$,
  show that the distance from $z$ to $F_1$ plus the distance from $z$ to $F_2$
  is a constant, so that the intersection is an ellipse, with foci $F_1$ and
  $F_2$.
\end{letters}

\begin{center}
\includesvg[width=0.62\textwidth,inkscapelatex=false]{assets/rendered/ch04-fig49}\\[0.6em]
\figlabel{8}
\end{center}

\pnum{3.} Similarly, use Figure~9(a) to prove geometrically that the
intersection of a plane and a cone is an ellipse when the plane intersects just
one half of the cone. Similarly, use~(b) to prove that the intersection is an
hyperbola when the plane intersects both halves of the cone.

\begin{center}
\includesvg[width=0.95\textwidth,inkscapelatex=false]{assets/rendered/ch04-fig50}\\[0.6em]
\figlabel{9}
\end{center}

\section*{APPENDIX 3.\quad POLAR COORDINATES}

In this chapter we've been acting all along as if there's only one way to label
points in the plane with pairs of numbers. Actually, there are many different
ways, each giving rise to a different ``coordinate system.'' The usual
coordinates of a point are called its cartesian coordinates, after the French
mathematician and philosopher René Descartes (1596--1650), who first introduced
the idea of coordinate systems. In many situations it is more convenient to
introduce polar coordinates, which are illustrated in Figure~1. To the point
$P$ we assign the polar coordinates $(r,\theta)$, where $r$ is the distance
from the origin $O$ to $P$, and $\theta$ is the measure, in radians, of the
angle between the horizontal axis and the line from $O$ to $P$. This $\theta$
is not determined unambiguously. For example, points on the right side of the
horizontal axis could have either $\theta=0$ or $\theta=2\pi$; moreover,
$\theta$ is completely ambiguous at the origin $O$. So it is necessary to
exclude some ray through the origin if we want to assign a unique pair
$(r,\theta)$ to each point under consideration.

\begin{center}
\figtex{assets/geometric/ch04-fig51}\\[0.6em]
\figlabel{1}
\end{center}

On the other hand, there is no problem associating a unique point to any pair
$(r,\theta)$. In fact, it is possible (though not approved of by all) to
associate a point to $(r,\theta)$ when $r<0$, according to the scheme indicated
in Figure~2. Thus, it always makes sense to talk about ``the point with polar
coordinates $(r,\theta)$,'' (with or without the possibility of $r<0$), even
though there is some ambiguity when we talk about ``the polar coordinates'' of
a given point.

\begin{center}
\figtex{assets/geometric/ch04-fig52}\\[0.6em]
\figlabel{2}
\end{center}

It is clear from Figure~1 (and Figure~2) that the point with polar coordinates
$(r,\theta)$ has cartesian coordinates $(x,y)$ given by
\[
x=r\cos\theta,\qquad y=r\sin\theta.
\]
Conversely, if a point has cartesian coordinates $(x,y)$, then (any of) its
polar coordinates $(r,\theta)$ satisfy
\begin{align*}
r&=\pm\sqrt{x^2+y^2}\\
\tan\theta&=\frac yx \qquad\text{if }x\neq 0.
\end{align*}

Now suppose that $f$ is a function. Then by the \textbf{graph of $f$ in polar
coordinates} we mean the collection of all points $P$ with polar coordinates
$(r,\theta)$ satisfying $r=f(\theta)$. In other words, the graph of $f$ in
polar coordinates is the collection of all points with polar coordinates
$(f(\theta),\theta)$. No special significance should be attached to the fact
that we are considering pairs $(f(\theta),\theta)$, with $f(\theta)$ first, as
opposed to pairs $(x,f(x))$ in the usual graph of $f$; it is purely a matter of
convention that $r$ is considered the first polar coordinate and $\theta$ is
considered the second.

The graph of $f$ in polar coordinates is often described as ``the graph of the
equation $r=f(\theta)$.'' For example, suppose that $f$ is a constant function,
$f(\theta)=a$ for all $\theta$. The graph of the equation $r=a$ is simply a
circle with center $O$ and radius $a$ (Figure~3). This example illustrates, in
a rather blatant way, that polar coordinates are likely to make things simpler
in situations that involve symmetry with respect to the origin $O$.

\begin{center}
\figtex{assets/geometric/ch04-fig53}\\[0.6em]
\figlabel{3}
\end{center}

The graph of the equation $r=\theta$ is shown in Figure~4. The solid line
corresponds to all values of $\theta\geq 0$, while the dashed line corresponds
to values of $\theta\leq 0$.

\begin{center}
\figtex{assets/geometric/ch04-fig54}\\[0.6em]
\figlabel{4}
\end{center}

As another example involving both positive and negative $r$, consider the graph
of the equation $r=\cos\theta$. Figure~5(a) shows the part that corresponds to
$0\leq\theta\leq\pi/2$. Figure~5(b) shows the part corresponding to
$\pi/2\leq\theta\leq\pi$; here $r<0$. You can check that no new points are added
for $\theta>\pi$ or $\theta<0$. It is easy to describe this same graph in terms
of the cartesian coordinates of its points. Since the polar coordinates of any
point on the graph satisfy
\[
r=\cos\theta,
\]
and hence
\[
r^2=r\cos\theta,
\]
its cartesian coordinates satisfy the equation
\[
x^2+y^2=x
\]
which describes a circle (Problem~4-16). [Conversely, it is clear that if the
cartesian coordinates of a point satisfy $x^2+y^2=x$, then it lies on the graph
of the equation $r=\cos\theta$.]

\begin{center}
\figtex{assets/geometric/ch04-fig55}\\[0.6em]
\figlabel{5}
\end{center}

Although we've now gotten a circle in two different ways, we might well be
hesitant about trying to find the equation of an ellipse in polar coordinates.
But it turns out that we can get a very nice equation if we choose one of the
\emph{foci} as the origin. Figure~6 shows an ellipse with one focus at $O$,
with the sum of the distances of all points from $O$ and the other focus
$\mathbf{f}$ being $2a$. We've chosen $\mathbf{f}$ to the left of $O$, with
coordinates written as
\[
(-2\varepsilon a,0).
\]
(We have $0\leq\varepsilon<1$, since we must have $2a>$ distance from
$\mathbf{f}$ to $O$).

\begin{center}
\figtex{assets/geometric/ch04-fig56}\\[0.6em]
\figlabel{6}
\end{center}

The distance $r$ from $(x,y)$ to $O$ is given by
\begin{equation}
r^2=x^2+y^2.\tag{1}
\end{equation}
By assumption, the distance from $(x,y)$ to $\mathbf{f}$ is $2a-r$, hence
\[
(2a-r)^2=\bigl(x-[-2\varepsilon a]\bigr)^2+y^2,
\]
or
\begin{equation}
4a^2-4ar+r^2=x^2+4\varepsilon ax+4\varepsilon^2a^2+y^2.\tag{2}
\end{equation}
Subtracting~(1) from~(2), and dividing by $4a$, we get
\[
a-r=\varepsilon x+\varepsilon^2 a,
\]
or
\begin{align*}
r&=a-\varepsilon x-\varepsilon^2 a\\
&=(1-\varepsilon^2)a-\varepsilon x,
\end{align*}
which we can write as
\begin{equation}
r=\Lambda-\varepsilon x, \qquad\text{for }\Lambda=(1-\varepsilon^2)a.\tag{3}
\end{equation}
Substituting $r\cos\theta$ for $x$, we have
\begin{align*}
r&=\Lambda-\varepsilon r\cos\theta,\\
r(1+\varepsilon\cos\theta)&=\Lambda,
\end{align*}
and thus
\begin{equation}
r=\frac{\Lambda}{1+\varepsilon\cos\theta}.\tag{4}
\end{equation}

In Chapter~4 we found that
\begin{equation}
\frac{x^2}{a^2}+\frac{y^2}{b^2}=1\tag{5}
\end{equation}
is the equation in cartesian coordinates for an ellipse with $2a$ as the sum of
the distances to the foci, but with the foci at $(-c,0)$ and $(c,0)$, where
\[
b=\sqrt{a^2-c^2}.
\]
Since the distance between the foci is $2c$, when this ellipse is moved left by
$c$ units, so that the focus $(c,0)$ is now at the origin, we get the
ellipse~(4) when we take $c=\varepsilon a$ or $\varepsilon=c/a$ (with
equation~(3) determining $\Lambda$). Conversely, given the ellipse described by
(4), for the corresponding equation~(5) the value of $a$ is determined by~(3),
\[
a=\frac{\Lambda}{1-\varepsilon^2},
\]
and again using $c=\varepsilon a$, we get
\[
b=\sqrt{a^2-c^2}=\sqrt{a^2-\varepsilon^2a^2}=a\sqrt{1-\varepsilon^2}
=\frac{\Lambda}{\sqrt{1-\varepsilon^2}}.
\]
Thus, we can obtain $a$ and $b$, the lengths of the major and minor axes,
immediately from $\varepsilon$ and $\Lambda$.

The number
\[
\varepsilon=\frac{c}{a}=\frac{\sqrt{a^2-b^2}}{a}=\sqrt{1-\Bigl(\frac{b}{a}\Bigr)^2},
\]
the \emph{eccentricity} of the ellipse, determines the ``shape'' of the ellipse
(the ratio of the major and minor axes), while the number $\Lambda$ determines
its ``size,'' as shown by~(4).

\subsection*{PROBLEMS}

\pnum{1.} If two points have polar coordinates $(r_1,\theta_1)$ and
$(r_2,\theta_2)$, show that the distance $d$ between them is given by
\[
d^2=r_1^2+r_2^2-2r_1r_2\cos(\theta_1-\theta_2).
\]
What does this say geometrically?

\pnum{2.} Describe the general features of the graph of $f$ in polar
coordinates if
\begin{romans}
\item $f$ is even.
\item $f$ is odd.
\item $f(\theta)=f(\theta+\pi)$.
\end{romans}

\pnum{3.} Sketch the graphs of the following equations.
\begin{romans}
\item $r=a\sin\theta$.
\item $r=a\sec\theta$. \textit{Hint:} It is a very simple graph!
\item $r=\cos 2\theta$. Good luck on this one!
\item $r=\cos 3\theta$.
\item $r=\abs{\cos 2\theta}$.
\item $r=\abs{\cos 3\theta}$.
\end{romans}

\pnum{4.} Find equations for the cartesian coordinates of points on the graphs
(i), (ii) and (iii) in Problem~3.

\pnum{5.} Consider a hyperbola, where the difference of the distance between
the two foci is the constant $2a$, and choose one focus at $O$ and the other at
$(-2\varepsilon a,0)$. (In this case, we must have $\varepsilon>1$). Show that
we obtain the exact same equation in polar coordinates
\[
r=\frac{\Lambda}{1+\varepsilon\cos\theta}
\]
as we obtained for an ellipse.

\pnum{6.} Consider the set of points $(x,y)$ such that the distance $(x,y)$ to
$O$ is equal to the distance from $(x,y)$ to the line $y=a$ (Figure~7). Show
that the distance to the line is $a-r\cos\theta$, and conclude that the
equation can be written
\[
a=r(1+\cos\theta).
\]
Notice that this equation for a parabola is again of the same form as~(4).

\begin{center}
\figtex{assets/geometric/ch04-fig57}\\[0.6em]
\figlabel{7}
\end{center}

\pnum{7.} Now, for any $\Lambda$ and $\varepsilon$, consider the graph in polar
coordinates of the equation~(4), which implies~(3). Show that the points
satisfying this equation satisfy
\[
(1-\varepsilon^2)x^2+y^2=\Lambda^2-2\Lambda\varepsilon x.
\]
Using Problem~4-16, show that this is an ellipse for $\varepsilon<1$, a
parabola for $\varepsilon=1$, and a hyperbola for $\varepsilon>1$.

\pnum{8.}
\begin{letters}
\item Sketch the graph of the \emph{cardioid} $r=1-\sin\theta$.
\item Show that it is also the graph of $r=-1-\sin\theta$.
\item Show that it can be described by the equation
  \[
  x^2+y^2=\sqrt{x^2+y^2}-y,
  \]
  and conclude that it can be described by the equation
  \[
  (x^2+y^2+y)^2=x^2+y^2.
  \]
\end{letters}

\pnum{9.} Sketch the graphs of the following equations.
\begin{romans}
\item $r=1-\frac12\sin\theta$.
\item $r=1-2\sin\theta$.
\item $r=2+\cos\theta$.
\end{romans}

\pnum{10.}
\begin{letters}
\item Sketch the graph of the \emph{lemniscate}
  \[
  r^2=2a^2\cos 2\theta.
  \]
\item Find an equation for its cartesian coordinates.
\item Show that it is the collection of all points $P$ in Figure~8 satisfying
  $d_1d_2=a^2$.
\item Make a guess about the shape of the curves formed by the set of all $P$
  satisfying $d_1d_2=b$, when $b>a^2$ and when $b<a^2$.
\end{letters}

\begin{center}
\figtex{assets/geometric/ch04-fig58}\\[0.6em]
\figlabel{8}
\end{center}
